\chapter{The Context Switch}
\label{ch:context-switch}

Every preemptive runtime turns on one small, brutal piece of code: the routine
that freezes one task's CPU state, thaws another's, and jumps. On a flat
register machine that routine is tedious but obvious. On the Xtensa LX7, with its
\emph{windowed} register file and a hardware FPU bolted on the side, it is neither
\textemdash{} and getting it wrong produces failures that look like anything but a
broken context switch. This chapter walks through the mechanism end to end: where
it lives, the state it must preserve, the two completely different ways it is
entered, and how the trickier parts were verified on hardware.

\section{Where it lives}\index{context switch}

The switch spans three files, two languages, and the boundary between the generic
GNARL kernel and this CPU port:

\begin{itemize}
  \item \texttt{crates/bb-runtimes/xtensa/context\_switch.S} \textemdash{} the
        register-level switch itself, in hand-written Xtensa assembly. It is
        copied verbatim into each profile's runtime tree
        (\texttt{<profile>-esp32s3/gnarl/context\_switch.S}) when
        \texttt{gen\_runtime.sh} builds the RTS, so editing the bb-runtimes copy
        and regenerating is what changes the running behaviour.
  \item \texttt{System.BB.CPU\_Primitives} (\texttt{s-bbcppr.adb}) \textemdash{}
        the Ada entry point \texttt{Context\_Switch}, a thin wrapper that picks
        the per-CPU thread-table slots and calls the assembly. This is the seam
        the portable kernel sees; every CPU port provides its own.
  \item \texttt{System.BB.CPU\_Specific} (\texttt{s-bbcpsp.ads}) \textemdash{} the
        \texttt{Context\_Buffer} record: the per-task saved CPU state, whose field
        offsets the assembly hard-codes.
\end{itemize}

Above all of this sits the portable scheduler, \texttt{System.BB.Threads.Queues}
(\texttt{s-bbthqu.adb}). It never touches a register; it only maintains the ready
queue and answers one question \textemdash{} \texttt{Context\_Switch\_Needed},
defined simply as \texttt{First\_Thread /= Running\_Thread}. When the highest
ready task is not the one on the CPU, the kernel calls
\texttt{CPU\_Primitives.Context\_Switch} and the machinery below takes over.

\section{The data structures}\index{context buffer}\index{stack pointer}\index{thread table}

Two structures carry the state.

\textbf{The \texttt{Context\_Buffer}}, one per task, is deliberately tiny:

\begin{longtable}{p{3cm} p{1.2cm} p{8cm}}
\toprule
\textbf{Field} & \textbf{Off} & \textbf{Meaning} \\ \midrule
\endhead
\texttt{SP}         & 0  & task stack pointer (top of its spilled frame) \\
\texttt{PC}         & 4  & resume target: a \texttt{retw} stub, or the new-task trampoline \\
\texttt{PS}         & 8  & processor state (INTLEVEL, window control, \dots) \\
\texttt{A0}         & 12 & windowed return address \\
\texttt{THREADPTR}  & 16 & thread-local-storage base \\
\texttt{CP\_State}  & 20 & pointer to the FPU/coprocessor save area \\
\texttt{SAR}        & 24 & shift-amount register \\
\texttt{LBEG/LEND/LCOUNT} & 28..36 & zero-overhead loop registers \\
\texttt{Frame\_Kind} & 40 & 0 = solicited, 1 = interrupt (see below) \\
\bottomrule
\end{longtable}

What is conspicuously \emph{absent} is the general register file. That is the
whole point of the windowed design, and the next section explains why.

\textbf{The thread tables}, \texttt{Running\_Thread\_Table} and
\texttt{First\_Thread\_Table}, are indexed by CPU (this is a genuine two-core SMP
target). \texttt{Running} is who is on the core now; \texttt{First} is who the
scheduler wants there. \texttt{Context\_Switch} passes the \emph{addresses} of
these two slots to the assembly, which dereferences them to the descriptors whose
first field is the \texttt{Context\_Buffer}. The switch updates
\texttt{Running := First} as part of resuming, so the table always names the live
task.

\section{The windowed-ABI premise}\index{register window}\index{window overflow}

On the Xtensa windowed ABI a \texttt{call} does not push registers; it
\emph{rotates} a 64-entry physical register file so the callee sees a fresh
window. Live data from outer frames stays in the physical registers until the
file fills, at which point the hardware raises window-overflow exceptions that
spill the oldest frames to their stacks; returning raises window-underflow
exceptions that reload them. The implication for a context switch is liberating:
to save a task, you do not copy sixteen or thirty-two registers \textemdash{} you
issue \texttt{SPILL\_ALL\_WINDOWS}, which flushes every live frame to that task's
own stack, and then almost the entire register state lives on the stack. All you
must record separately is the stack pointer and the handful of special registers
the window machinery does not cover (\texttt{SAR}, the loop registers, the FPU,
\texttt{THREADPTR}). Resuming is the mirror image: set the stack pointer, execute
one \texttt{retw}, and let window-underflow lazily fault the frames back in as the
task returns up its own call chain. This is why the \texttt{Context\_Buffer} is so
small, and why this port is closest in spirit to the SPARC (LEON) windowed port
rather than the flat ARM and RISC-V ones.

\section{Two ways in}\index{preemption}\index{interrupt frame}

A task can lose the CPU two very different ways, and the switch handles them with
two different entry points.

\textbf{Cooperative} \textemdash{} the running task blocks (a \texttt{delay
until}, a protected entry that cannot proceed) or a higher task becomes ready at
a point where the kernel is already running. The kernel calls
\texttt{Context\_Switch}, which calls \texttt{\_\_gnat\_context\_switch}. This path
runs in the task's own context, so it can do the full \texttt{SPILL\_ALL\_WINDOWS}
and save the outgoing task in \emph{solicited} form.

\textbf{Preemptive} \textemdash{} a level-5 timer or IPI interrupt fires while a
task is running. The L5 vector (\texttt{highint5.S}) saves the interruptee with
\texttt{\_xt\_context\_save} into a full \texttt{XT\_STK} interrupt frame, services
the tick, and \textemdash{} if that made a higher task ready \textemdash{}
tail-jumps \texttt{\_\_gnat\_preempt\_dispatch} to switch away. This path saves the
outgoing task in \emph{interrupt} form.

The two must not be mixed, and that is the subtle part. A cooperative
\texttt{SPILL\_ALL\_WINDOWS} executed from inside an interrupt's borrowed window
chain corrupts the register windows. So when the kernel decides mid-interrupt that
a switch is needed, \texttt{Context\_Switch} does \emph{not} switch: it sets a
per-CPU \texttt{Switch\_Pending} flag and returns. The interrupt epilogue, running
in the vector's clean single-window context, notices the flag and performs the
dispatch itself via \texttt{\_\_gnat\_preempt\_dispatch}. (Tracking down a bug here
\textemdash{} the cooperative spill running in interrupt context, marching the
stack pointer down one frame per tick until a window spill double-faulted
\textemdash{} is the story behind several of the comments in \texttt{highint5.S}
and \texttt{context\_switch.S}, and behind ACATS \texttt{CXD8002}.)

\section{Dual-format dispatch}\index{WINDOWBASE}\index{PS register}

Because a task can be saved either way, each \texttt{Context\_Buffer} records its
\texttt{Frame\_Kind}, and the shared resume code reads it to resume correctly:

\begin{itemize}
  \item \textbf{Solicited} (\texttt{Frame\_Kind = 0}): the live windows were
        spilled to the task's stack. Resume by restoring \texttt{SP}, \texttt{SAR},
        the loop registers and \texttt{PS}, then a single \texttt{retw} \textemdash{}
        window-underflow reloads the rest.
  \item \textbf{Interrupt} (\texttt{Frame\_Kind = 1}): the full register state is in
        an \texttt{XT\_STK} frame at \texttt{SP}. Resume by calling
        \texttt{\_xt\_context\_restore} and returning with \texttt{rfi}.
\end{itemize}

Both entry points converge on a single label, \texttt{.Lresume}: it makes
\texttt{First} the running task, restores the format-independent extra state (FPU,
\texttt{THREADPTR}), then branches on \texttt{Frame\_Kind} to one of the two
resume tails. The interrupt tail carries hard-won detail \textemdash{} it forces a
single clean window (\texttt{WINDOWSTART = 1 << WINDOWBASE}) so a freshly scheduled
task does not inherit the dispatcher's stale window mask, reloads \texttt{SAR}
(clobbered by the window-start arithmetic), and returns with \texttt{rfi 5} so
\texttt{INTLEVEL} stays high until the very last instruction, closing the race
where a fresh tick preempts a half-restored task.

\section{Walkthrough: a cooperative switch}

A task calls \texttt{delay until}. The kernel marks it blocked, finds a different
\texttt{First\_Thread}, and calls \texttt{\_\_gnat\_context\_switch} with the two
slot addresses:

\begin{enumerate}
  \item \texttt{SPILL\_ALL\_WINDOWS} flushes the task's live frames to its stack.
  \item Save into the outgoing buffer: \texttt{SP}, \texttt{A0}, \texttt{PS},
        \texttt{SAR}, the loop registers, \texttt{THREADPTR}, and (if the task owns
        an FPU save area) \texttt{F0..F15} + \texttt{FCR} + \texttt{FSR}.
        \texttt{Frame\_Kind := 0}.
  \item Fall into \texttt{.Lresume}: set \texttt{Running := First}, restore the
        incoming task's FPU and \texttt{THREADPTR}, and \textemdash{} since the
        incoming task was itself last saved solicited \textemdash{} restore its
        \texttt{SP}/\texttt{SAR}/loops/\texttt{PS} and \texttt{retw} into it.
\end{enumerate}

\section{Walkthrough: a preemptive switch}

The L5 timer fires while a compute task runs:

\begin{enumerate}
  \item The vector saves the interruptee with \texttt{\_xt\_context\_save} into an
        \texttt{XT\_STK} frame and services the tick.
  \item If a higher task is now ready, the epilogue tail-jumps
        \texttt{\_\_gnat\_preempt\_dispatch} from the vector's clean single window
        (no spill here \textemdash{} the full state is already in \texttt{XT\_STK}).
  \item \texttt{\_\_gnat\_preempt\_dispatch} records the interruptee:
        \texttt{SP} (= the \texttt{XT\_STK} frame), \texttt{THREADPTR}, the FPU,
        \texttt{Frame\_Kind := 1}.
  \item \texttt{.Lresume} restores the incoming task's FPU/\texttt{THREADPTR} and
        resumes it in its own format \textemdash{} solicited via \texttt{retw}, or
        interrupt via \texttt{\_xt\_context\_restore} + \texttt{rfi}.
\end{enumerate}

\section{Starting a task that has never run}\index{task stack}\index{trampoline}

A brand-new task has no spilled frames to reload, so its first dispatch cannot use
\texttt{retw}. \texttt{Initialize\_Context} sets the new buffer's \texttt{PC} field
to the trampoline \texttt{\_\_gnat\_start\_thread} and stashes the task's entry
point and argument at the top of its fresh stack. The first time the scheduler
resumes the task, control lands in the trampoline, which loads the argument and
\texttt{callx} es the task wrapper \textemdash{} which never returns. The buffer's
loop registers are zero-initialised (a non-zero \texttt{LCOUNT} would make the
hardware spuriously loop on first entry) and \texttt{THREADPTR} starts at zero,
i.e.\ ``no thread-local storage'' until the task sets it.

\section{What state is saved \textemdash{} and what is not}\index{coprocessor state}\index{THREADPTR}

\begin{longtable}{p{5.5cm} p{2.2cm} p{5.5cm}}
\toprule
\textbf{State} & \textbf{Saved?} & \textbf{Notes} \\ \midrule
\endhead
\texttt{SP}/\texttt{PS}/\texttt{A0}, register windows & always & spill + frame \\
\texttt{SAR}, \texttt{LBEG}/\texttt{LEND}/\texttt{LCOUNT} & always & loops: \texttt{XCHAL\_HAVE\_LOOPS}=1 \\
\texttt{THREADPTR} & always & TLS base; per task \\
\texttt{F0..F15} + \texttt{FCR} + \texttt{FSR} & conditional & only when the task owns a \texttt{CP\_State} area \\
MAC16 \texttt{ACCLO}/\texttt{ACCHI}/\texttt{Mn} & no & \texttt{XCHAL\_HAVE\_MAC16}=0 \textemdash{} the unit is absent \\
\bottomrule
\end{longtable}

The FPU is handled \emph{eagerly per task}: the run-time enables the FPU bit
(\texttt{CP0}) of \texttt{CPENABLE} for every task, gives each a \texttt{CP\_State}
save area, and the switch copies all sixteen \texttt{F} registers plus
\texttt{FCR}/\texttt{FSR} whenever the buffer's \texttt{CP\_State} pointer is
non-null. Note this covers \emph{only} the FPU: the PIE/SIMD coprocessor
(\texttt{CP3}) is not part of the per-task enable, so a task using
\texttt{ESP32S3.SIMD} must switch it on itself (\S\ref{ch:asm}) and keep its vector
work single-task, since the \texttt{q}-registers are not in the save area. \texttt{THREADPTR} was the last addition to
the list; for some time the field existed in the record but the assembly never
touched it, leaving all tasks sharing one value \textemdash{} harmless only because
nothing yet used per-task TLS.

\section{Verifying the hard parts}

The FPU and \texttt{THREADPTR} paths share a verification hazard: a test that
exercises them na\"ively passes whether or not the switch actually preserves them,
because the value under test was sitting safely in memory the whole time. A green
test means nothing unless it is \emph{sensitive} \textemdash{} demonstrably red
when the property is broken. Both were therefore verified two-sided on the
\texttt{full} profile: a deliberately broken build must fail, and only then does
the intact build's success count.

\textbf{FPU.} A low-priority checker holds six floating-point accumulators live in
registers across a two-million-step loop (each step the exact identity
\texttt{X := X * Lm * Li}, with \texttt{Lm}=2.0 and \texttt{Li}=0.5 read once from
\texttt{Volatile} cells so the compiler can neither fold the pair away nor spill
the accumulators \textemdash{} disassembly confirmed zero memory operations in the
loop). A higher-priority hog wakes every millisecond and tramples all of
\texttt{F0..F15}. With the preemptive FPU save disabled the checker logs corruption
continuously; with it intact, zero corruption across hundreds of preemptions per
batch.

\textbf{THREADPTR.} A checker task sets \texttt{THREADPTR} to one sentinel; a
higher-priority hog sets a different one and preempts it (both confirmed on the
same core). The clean discriminator turned out to be the \emph{hog's} value on
re-entry, not the checker's \textemdash{} the checker, as the background task,
reads its own value either way. With the restore disabled the hog comes back
seeing the checker's sentinel (its own was lost); with the restore intact it keeps
its own. That asymmetry is the proof the per-task save/restore works.

Two methodology traps recurred and are worth stating once. First, optimisation:
a value only tests the switch if the compiler keeps it \emph{live in a register}
across the preemption window \textemdash{} verify that by disassembling the loop,
not by reading the source. Second, staleness: editing
\texttt{bb-runtimes/xtensa/context\_switch.S} does not take effect until the
generated \emph{base} runtime tree is deleted and regenerated, and the only
trustworthy confirmation that a change is in the running image is to disassemble
the final linked ELF. Both traps produce a green run from a binary that does not
contain your change \textemdash{} the most dangerous kind of pass.
